\chapter{Data Quality and Exploratory Analysis}
\label{ch:eda}

\section{Theory}

\subsection{Data quality auditing}
Before any modelling, three questions must be answered about a tabular
dataset $D = \{(\mathbf{x}_i)\}_{i=1}^n$: (i) is every field typed and
parsed correctly, (ii) is any value missing, and (iii) is any row
duplicated. Formally, for a column $X$ we report the \emph{missingness
rate} $\hat\pi_X = \frac{1}{n}\sum_{i=1}^n \mathbb{1}[x_{i} \text{ is
NA}]$ and \emph{cardinality} $|\{x_i\}|$. A duplicate row under a primary
key $k$ is any pair $i \neq j$ with $k_i = k_j$; if such pairs exist, the
key does not uniquely identify a record and downstream joins or grouped
aggregates will silently double-count. These checks are cheap and
guard against exactly the class of error (a re-exported CSV with
duplicated header rows, an encoding mismatch producing NA in a text field)
that otherwise corrupts every later stage silently.

\subsection{Descriptive statistics and distribution shape}
For a numeric column $X$ with observed values $x_1,\dots,x_n$, the sample
mean $\bar x = \frac1n\sum_i x_i$ and variance
$s^2 = \frac{1}{n-1}\sum_i (x_i-\bar x)^2$ summarise location and spread,
but retail money variables are rarely well described by them alone because
they are \emph{right-skewed}: many small transactions and a long tail of
large ones. Skewness is measured by the third standardised moment,
\begin{equation}
  \gamma_1 = \frac{\frac1n\sum_i (x_i - \bar x)^3}{\left(\frac1n\sum_i (x_i-\bar x)^2\right)^{3/2}},
  \label{eq:skew}
\end{equation}
with $\gamma_1 > 0$ indicating a right tail. A common remedy that
stabilises both skew and the variance-vs-mean relationship
(heteroscedasticity) seen in multiplicative processes like sales is the
$\log(1+x)$ transform, used throughout this project on \texttt{Sales}
(and, with a sign-preserving variant, on \texttt{Profit}) purely for
\emph{visualisation}; models are always trained on the original scale
unless stated otherwise, since tree ensembles are invariant to monotone
transforms and linear models' coefficients are more directly interpretable
untransformed here.

\subsection{Correlation}
For two numeric variables $X, Y$, the Pearson correlation coefficient
\begin{equation}
  \rho_{XY} = \frac{\mathrm{Cov}(X,Y)}{\sigma_X \sigma_Y}
            = \frac{\sum_i (x_i - \bar x)(y_i - \bar y)}
                   {\sqrt{\sum_i (x_i-\bar x)^2}\sqrt{\sum_i (y_i-\bar y)^2}}
  \label{eq:pearson}
\end{equation}
measures \emph{linear} association, $\rho_{XY}\in[-1,1]$. It is
insensitive to non-linear relationships (e.g. a U-shape can have
$\rho \approx 0$ despite strong dependence), which is precisely why
\cref{sec:eda-discount-profit} supplements the correlation matrix with a
binned conditional-mean plot: discount and profit turn out to have a
threshold-like, not linear, relationship.

\subsection{Seasonal-trend decomposition (STL)}
A time series $y_t$ is classically modelled as the additive sum of a
slowly-varying trend $T_t$, a periodic seasonal component $S_t$ (period
$L$, here $L=12$ months), and a remainder $R_t$:
\begin{equation}
  y_t = T_t + S_t + R_t.
  \label{eq:stl-additive}
\end{equation}
\textbf{STL} (Seasonal-Trend decomposition using \textsc{Loess}, Cleveland
et al., 1990) estimates $T_t$ and $S_t$ by alternating two local
regression (LOESS) smooths: it removes a running estimate of $T_t$ from
$y_t$ to expose the seasonal cycle, smooths within each seasonal
sub-series (all Januaries, all Februaries, \dots) to obtain a smoothly
time-varying $S_t$, then smooths $y_t - S_t$ to update $T_t$; this is
iterated to convergence. Unlike a single global seasonal index, STL lets
the seasonal shape drift slowly over the four years of data, and unlike
classical decomposition it is robust to outliers when the \texttt{robust}
option down-weights large residuals between iterations. We use STL purely
as a diagnostic here; the forecasting models of \cref{ch:forecasting} are
fitted directly on $y_t$, not on the decomposed components.

\section{Method}
The raw CSV is read with \texttt{src.data\_loader.load\_raw} (encoding
auto-detection across UTF-8/Latin-1/CP1252) and typed with \texttt{clean},
which parses \texttt{Order Date}/\texttt{Ship Date}, strips whitespace from
every text column, drops exact and \texttt{Row ID} duplicates, and derives
\texttt{Ship Days} $= \text{Ship Date} - \text{Order Date}$. The quality
audit, univariate/bivariate plots and STL decomposition are produced by
notebooks \texttt{01\_data\_overview.ipynb} and \texttt{02\_eda.ipynb}.

\section{Results}

\subsection{Data quality}
The audit (\texttt{outputs/tables/01\_data\_quality.csv}) finds
\textbf{zero missing values in any of the 21 raw columns}, no fully
duplicated rows, and no duplicate \texttt{Row ID}s after cleaning — the
data is analysis-ready without imputation. \texttt{Order Date} spans
2014-01-03 to 2017-12-30 (1{,}237 distinct dates); \texttt{Ship Days} is
always non-negative, i.e. no order ships before it is placed.

\subsection{Distributions}
\begin{figure}[H]
  \centering
  \includegraphics[width=0.95\linewidth]{01_raw_distributions.png}
  \caption{Histograms (log-count $y$-axis) of \texttt{Sales}, \texttt{Profit}
  and \texttt{Discount} on the raw scale, showing pronounced right skew in
  \texttt{Sales} and a bimodal, discount-tier structure in \texttt{Discount}.}
  \label{fig:raw-dist}
\end{figure}

\begin{figure}[H]
  \centering
  \includegraphics[width=0.95\linewidth]{02_log_transforms.png}
  \caption{Left: $\log(1+\text{Sales})$ recovers an approximately unimodal,
  near-symmetric shape (\cref{eq:skew} in the raw scale is large and
  positive; after the transform it is close to zero). Right: signed
  $\log(1+|\text{Profit}|)$, showing the loss-making mass as a distinct
  negative lobe rather than a long left tail.}
  \label{fig:log-dist}
\end{figure}

\begin{figure}[H]
  \centering
  \includegraphics[width=0.95\linewidth]{02_univariate_categorical.png}
  \caption{Order-line counts by \texttt{Ship Mode}, \texttt{Segment},
  \texttt{Region} and \texttt{Category}. Standard Class and the Consumer
  segment dominate; the four regions are comparably sized, with West
  largest.}
  \label{fig:cat-dist}
\end{figure}

\subsection{Profit by product hierarchy}
Table~\ref{tab:subcat-profit} (full detail in
\texttt{outputs/tables/02\_by\_subcategory.csv}) aggregates profit by
sub-category. \textbf{Tables} lose \$17{,}725 overall despite \$207k of
sales, and \textbf{Bookcases} lose \$3{,}473 — both at average discounts
above 21\%, well above the store-wide mean. \textbf{Supplies} is also
marginally loss-making. Every other sub-category is net profitable, led by
Copiers (37\% margin) and Phones (\$330k sales, 13\% margin).

\begin{table}[H]
  \centering
  \caption{Selected sub-categories by total profit (ascending). Full table
  has all 17 sub-categories.}
  \label{tab:subcat-profit}
  \begin{tabular}{@{}llrrrr@{}}
    \toprule
    Category & Sub-Category & Sales (\$) & Profit (\$) & Avg.\ Discount & Margin \\
    \midrule
    Furniture & Tables      & 206{,}966 & $-$17{,}725 & 26.1\% & $-$8.6\% \\
    Furniture & Bookcases   & 114{,}880 & $-$3{,}473  & 21.1\% & $-$3.0\% \\
    Office Supplies & Supplies & 46{,}674 & $-$1{,}189 & 7.7\% & $-$2.5\% \\
    \multicolumn{6}{c}{\dots (14 profitable sub-categories) \dots} \\
    Technology & Phones     & 330{,}007 & 44{,}516    & 15.5\% & 13.5\% \\
    Technology & Copiers    & 149{,}528 & 55{,}618    & 16.2\% & \textbf{37.2\%} \\
    \bottomrule
  \end{tabular}
\end{table}

\begin{figure}[H]
  \centering
  \includegraphics[width=0.8\linewidth]{02_profit_by_subcategory.png}
  \caption{Total profit by sub-category; red bars are net loss-makers.}
  \label{fig:subcat-profit}
\end{figure}

\begin{figure}[H]
  \centering
  \includegraphics[width=0.85\linewidth]{02_region_segment_heatmap.png}
  \caption{Sales (left) and profit (right) cross-tabulated by Region and
  Segment. Sales are fairly even across cells; profit is not — the
  Central/Consumer cell is comparatively weak.}
  \label{fig:region-segment}
\end{figure}

\subsection{Discount versus profit}
\label{sec:eda-discount-profit}
\begin{figure}[H]
  \centering
  \includegraphics[width=0.95\linewidth]{02_discount_vs_profit.png}
  \caption{Left: line-level profit vs.\ discount (3,000-row sample),
  coloured by category — profit is scattered near zero at low discount and
  swings sharply negative as discount grows. Right: mean profit by
  discount band, crossing zero between the 10--20\% and 20--30\% bands and
  falling steeply beyond 30\%.}
  \label{fig:discount-profit}
\end{figure}
Quantitatively, the loss rate ($\Pr(\text{Profit}<0)$) is
\textbf{exactly 0\%} for undiscounted lines and \textbf{96.8\%} for lines
discounted 30\% or more, against an overall loss rate of 18.7\% across all
9{,}994 lines — discount depth is close to a deterministic switch, not a
mild risk factor. Because the relationship is a threshold
effect rather than linear, the raw Pearson correlation between
\texttt{Discount} and \texttt{Profit} (visible in
\cref{fig:corr-heatmap}) understates its practical importance; this
motivates using the categorical \texttt{discount\_bucket} feature (rather
than relying on a linear term alone) in \cref{ch:features,ch:models}.

\begin{figure}[H]
  \centering
  \includegraphics[width=0.55\linewidth]{02_correlation.png}
  \caption{Pearson correlation (\cref{eq:pearson}) among the core numeric
  fields. \texttt{Sales}--\texttt{Profit}: $\rho\approx0.48$;
  \texttt{Discount}--\texttt{Profit}: $\rho\approx-0.22$ (negative but,
  per \cref{fig:discount-profit}, understating the non-linear effect);
  \texttt{Quantity} is nearly uncorrelated with profit.}
  \label{fig:corr-heatmap}
\end{figure}

\subsection{Seasonality}
\begin{figure}[H]
  \centering
  \includegraphics[width=0.95\linewidth]{02_stl_decomposition.png}
  \caption{STL decomposition (\cref{eq:stl-additive}, period 12) of monthly
  sales: a clear upward trend across the four years and a recurring
  within-year seasonal pattern that peaks late in the year, leaving a
  comparatively small, structure-free remainder.}
  \label{fig:stl}
\end{figure}
\begin{figure}[H]
  \centering
  \includegraphics[width=0.95\linewidth]{02_seasonal_profiles.png}
  \caption{Sales aggregated by month-of-year (left) and day-of-week
  (right, 0 = Monday). November--December is by far the strongest period;
  day-of-week shows only mild variation.}
  \label{fig:seasonal-profiles}
\end{figure}

\subsection{Customers and products}
The top-10 customers by sales (Table~\ref{tab:top-cust}) generate between
\$12k and \$25k each, but profitability is uneven within that list —
\emph{Sean Miller} is the single largest customer by revenue (\$25{,}043)
yet is net loss-making ($-\$1{,}981$), while \emph{Tamara Chand} generates
comparable revenue at a much healthier margin. The same pattern repeats at
product level: the \emph{Cisco TelePresence System EX90} and two
\emph{GBC} binding systems are top-10 by revenue but net loss-making,
again tracing back to discount depth on high-ticket Technology and
Furniture items.

\begin{table}[H]
  \centering
  \caption{Top 5 customers by total sales.}
  \label{tab:top-cust}
  \begin{tabular}{@{}lrrr@{}}
    \toprule
    Customer & Sales (\$) & Profit (\$) & Orders \\
    \midrule
    Sean Miller     & 25{,}043 & $-$1{,}981 & 5 \\
    Tamara Chand    & 19{,}052 & 8{,}981    & 5 \\
    Raymond Buch    & 15{,}117 & 6{,}976    & 6 \\
    Tom Ashbrook    & 14{,}596 & 4{,}704    & 4 \\
    Adrian Barton   & 14{,}474 & 5{,}445    & 10 \\
    \bottomrule
  \end{tabular}
\end{table}

\section{Interpretation}
Three findings from this chapter shape every subsequent stage of the
report. First, the data requires \emph{no} missing-value or duplicate
handling, so all modelling error downstream is genuinely about the
signal, not data-cleaning artefacts. Second, \textbf{discount depth, not
category or region, is the primary lever on profitability} — a threshold
effect around 30\% discount that is far stronger than any single-variable
linear correlation suggests, which is why \texttt{discount\_bucket}
(\cref{ch:features}) and the corresponding permutation-importance results
(\cref{ch:models}) recur as the dominant signal throughout the rest of the
report. Third, the pronounced trend and December seasonality in
\cref{fig:stl} is exactly the structure a seasonal model must capture,
and directly motivates the choice of Holt-Winters and SARIMA — both
explicitly seasonal — over simpler alternatives in \cref{ch:forecasting}.
